% ------------------------------------------------------------------------------
\section{Modelo 3D, metrado e intercambio}
\label{sec:mod}
% ------------------------------------------------------------------------------

Este capítulo describe cómo ConcretoLab representa la geometría del edificio, cómo la genera (a mano, con plantillas,
por ejes y niveles o desde un plano CAD), cómo la edita y la mueve, cómo calcula el metrado sin doble conteo y cómo la
intercambia con otros programas (DXF, IFC y ETABS). Todo lo que se afirma aquí se obtuvo del código fuente; cuando un
valor depende de una norma o de un formato externo y no se ha podido comprobar, se indica expresamente.

% ------------------------------------------------------------------------------
\subsection{Modelo geométrico y proyectos}
\label{sec:mod-modelo}

\subsubsection{La entidad \texorpdfstring{\ruta{Elemento}}{Elemento}}

Cada pieza del edificio es un \textbf{prisma recto alineado con los ejes globales} (entidad JPA \ruta{model/Elemento},
tabla \ruta{elementos}). Los ejes X e Y son horizontales y Z es vertical. Las coordenadas $(x, y, z)$ guardadas son las
del \textbf{centro} del prisma, y las tres medidas se miden a lo largo de los ejes, no de la dirección «natural» del
elemento.

\begin{table}[H]
	\centering
	\caption{Campos de \ruta{Elemento}.}
	\label{tab:mod-elemento}
	\begin{tabular}{L{2.6cm}L{2.4cm}L{8.6cm}}
		\hline
		\textbf{Campo} & \textbf{Tipo} & \textbf{Significado} \\ \hline
		\ruta{id} & \ruta{Long} & Identificador autonumérico. \\
		\ruta{proyectoId} & \ruta{Long} & Proyecto al que pertenece (obligatorio). \\
		\ruta{tipo} & \ruta{String} & \ruta{ZAPATA}, \ruta{COLUMNA}, \ruta{VIGA}, \ruta{LOSA} o \ruta{MURO}. \\
		\ruta{nombre} & \ruta{String} & Marca del elemento (obligatoria, no se exige que sea única). \\
		\ruta{largo} & \ruta{double} & Medida a lo largo de X (m). \\
		\ruta{ancho} & \ruta{double} & Medida a lo largo de Y (m). \\
		\ruta{altura} & \ruta{double} & Medida a lo largo de Z (m). \\
		\ruta{x}, \ruta{y}, \ruta{z} & \ruta{double} & Coordenadas del centro (m). \\
		\hline
	\end{tabular}
\end{table}

De esta convención se derivan los límites del prisma y su volumen bruto, que devuelve \ruta{Elemento.volumen()}:

\begin{align}
	x^{\pm} &= x \pm \tfrac{l}{2}, & y^{\pm} &= y \pm \tfrac{a}{2}, & z^{\pm} &= z \pm \tfrac{h}{2},
	\label{eq:mod-limites} \\
	V_{\text{bruto}} &= l \, a \, h ,
\end{align}

\noindent donde $l$ es el largo, $a$ el ancho y $h$ la altura. Por lo tanto, una viga en dirección X tiene su luz en
\ruta{largo}, mientras que una viga en dirección Y la tiene en \ruta{ancho}. Varias partes del código deducen la
dirección de una viga comparando ambas medidas (véase la sección~\ref{sec:mod-limites}).

\begin{importante}[Medidas según los ejes globales]
	Al editar una viga, recuerde que «Largo eje X» es su longitud solo si la viga va en X. El aviso de vigas sueltas
	(sección~\ref{sec:mod-sueltas}) lo recuerda en su mensaje, porque el error más frecuente es fijar la misma medida a
	vigas de ambas direcciones.
\end{importante}

\subsubsection{Validación al guardar}

\ruta{YefriLermaProyectoService.\allowbreak{}yefriLermaGuardar(Elemento)} rechaza el elemento si el proyecto no existe, si el tipo
no es uno de los cinco admitidos, si el nombre está vacío, si alguna medida no es un número finito mayor que cero, si el
volumen no es finito o si alguna coordenada no es finita. Todas las vías de creación (formulario, plantilla, cuadrícula
e importación DXF) terminan llamando a este método.

\ruta{yefriLermaGuardarVarios(List<Elemento>)} guarda una lista dentro de una única transacción (\ruta{@Transactional}):
si un elemento no es válido, la excepción revierte la transacción y \textbf{no se guarda ninguno}. Es la base de la
edición en lote y del deshacer (sección~\ref{sec:mod-pantalla}).

\subsubsection{Proyectos}

\begin{itemize}
	\item \ruta{yefriLermaCrearProyecto(nombre, ubicacion)} exige un nombre de 1 a 120 caracteres (tras recortar
	      espacios) y una ubicación de hasta 200, y registra \ruta{PROYECTO\_CREADO} en la bitácora.
	\item \ruta{yefriLermaEliminarProyecto(id)} borra, en este orden, los ensayos de cada mezcla, las mezclas, los
	      elementos (cada uno con sus vínculos viga--columna, mediante \ruta{yefriLermaEliminar}), la cuadrícula, el
	      expediente y el proyecto, y registra \ruta{PROYECTO\_ELIMINADO}.
	\item \ruta{yefriLermaGuardarPlano(id, nombre, candidatos)} guarda los elementos importados de un plano CAD en un
	      proyecto nuevo (si \ruta{id} es nulo) o en uno existente, que debe estar \textbf{vacío}.
\end{itemize}

\subsubsection{Plantilla rectangular}

\ruta{yefriLermaCrearPlantilla(proyectoId, Plantilla)} crea un pórtico de un vano por dirección: cuatro zapatas, cuatro
columnas y cuatro vigas. Con separaciones $S_x$, $S_y$, zapata $l_z \times a_z \times h_z$, columna cuadrada de lado
$b_c$ y altura $H_c$, y viga $b_v \times h_v$, las posiciones son:

\begin{align}
	x_i &\in \left\{-\tfrac{S_x}{2},\, \tfrac{S_x}{2}\right\}, \qquad y_j \in \left\{-\tfrac{S_y}{2},\, \tfrac{S_y}{2}\right\}, \\
	z_{\text{zapata}} &= \tfrac{h_z}{2}, \qquad z_{\text{columna}} = h_z + \tfrac{H_c}{2}, \qquad
	z_{\text{viga}} = h_z + H_c - \tfrac{h_v}{2}, \\
	\text{VX}_{1,2} &: \; (S_x - b_c) \times b_v \times h_v \text{ con centro en } (0,\, y_j,\, z_{\text{viga}}), \\
	\text{VY}_{1,2} &: \; b_v \times (S_y - b_c) \times h_v \text{ con centro en } (x_i,\, 0,\, z_{\text{viga}}).
\end{align}

La base de las zapatas queda en $z = 0$. El método exige todas las medidas positivas, $S_x > b_c$, $S_y > b_c$,
$h_v \le H_c$ y $b_v \le b_c$.

\subsubsection{Escenarios de laboratorio}

\ruta{yefriLermaPrepararEscenarios()} crea, si no existen ya con ese nombre, los tres proyectos de
\ruta{YEFRI\_LERMA\_ESCENARIOS}, pensados para comparar respuestas dinámicas. Todos usan una planta de $3 \times 2$ ejes
con $x \in \{-5, 0, 5\}$\,m e $y \in \{-2, 2\}$\,m, zapatas de $2 \times 2 \times 0{,}65$\,m solo en el primer nivel,
columnas que arrancan en $z = 0{,}65$\,m y vigas de $0{,}30 \times 0{,}50$\,m.

\begin{table}[H]
	\centering
	\caption{Escenarios de laboratorio que crea \ruta{yefriLermaPrepararEscenarios}.}
	\label{tab:mod-escenarios}
	\begin{tabular}{L{5.2cm}cL{6.6cm}}
		\hline
		\textbf{Nombre} & \textbf{Pisos} & \textbf{Altura y columna por piso} \\ \hline
		LAB, 4 pisos, pórtico regular & 4 & 3,10\,m; columnas de 0,45\,m \\
		LAB, 3 pisos, planta baja flexible & 3 & piso 1: 4,50\,m y columnas de 0,30\,m; demás: 3,10\,m y 0,45\,m \\
		LAB, 3 pisos, referencia regular & 3 & 3,10\,m; columnas de 0,45\,m \\
		\hline
	\end{tabular}
\end{table}

\begin{nota}[Sin llamada desde la interfaz]
	En el código revisado, \ruta{yefriLermaPrepararEscenarios} solo se invoca desde la prueba
	\ruta{ProyectoIntegracionTest}; ninguna clase de \ruta{src/main} la llama. Los nombres reales usan un punto medio
	como separador («LAB · 4 pisos · pórtico regular»).
\end{nota}

% ------------------------------------------------------------------------------
\subsection{Ejes, niveles y generación}
\label{sec:mod-cuadricula}

El diálogo «Ejes, niveles y generación» (\ruta{ui/EditorCuadricula}) y el servicio \ruta{service/CuadriculaService}
generan zapatas, columnas y vigas sobre una cuadrícula. La definición (\ruta{Definicion}) tiene un origen
$(x_0, y_0)$, una lista de separaciones en X, otra en Y y una lista de niveles $N_0 < N_1 < \dots < N_m$. Las
coordenadas de los ejes salen de \ruta{walterQuispesayhuaCoordenadas}:

\begin{equation}
	X_0 = x_0, \qquad X_i = X_{i-1} + s^{x}_{i} \quad (i = 1, \dots, n_x - 1),
\end{equation}

\noindent y análogamente para $Y_j$. El primer nivel $N_0$ es la cara superior de las zapatas.

\begin{table}[H]
	\centering
	\caption{Validaciones de \ruta{walterQuispesayhuaValidarDefinicion} y de la generación.}
	\label{tab:mod-cuadricula-limites}
	\begin{tabular}{L{5.2cm}L{8.4cm}}
		\hline
		\textbf{Dato} & \textbf{Regla} \\ \hline
		Origen & Números finitos. \\
		Separaciones X e Y & De 1 a 19 valores por dirección (2 a 20 ejes), todos positivos y finitos. \\
		Niveles & De 2 a 20 valores, finitos y estrictamente crecientes. \\
		Medidas de zapata, columna y viga & Positivas y finitas. \\
		Ancho de viga & $b_v \le b_c$ si se generan vigas. \\
		Peralte de viga & $h_v \le N_k - N_{k-1}$ en cada piso si se generan vigas. \\
		Luz libre & $X_{i+1} - X_i - b_c > 0$ (y lo mismo en Y). \\
		Cantidad & Como máximo 2000 elementos por operación. \\
		\hline
	\end{tabular}
\end{table}

\subsubsection{Elementos generados}

Para cada intersección $(i, j)$ seleccionada, con marca «$i{+}1$-$j{+}1$», \ruta{walterQuispesayhuaGenerar} produce:

\begin{align}
	\text{Zapata} &: \; l_z \times a_z \times h_z, \quad (X_i,\, Y_j,\, N_0 - \tfrac{h_z}{2}), \\
	\text{Columna del piso } k &: \; b_c \times b_c \times H_k, \quad (X_i,\, Y_j,\, N_{k-1} + \tfrac{H_k}{2}),
	\qquad H_k = N_k - N_{k-1}, \\
	\text{Viga X del piso } k &: \; (X_{i+1} - X_i - b_c) \times b_v \times h_v, \quad
	\left(\tfrac{X_i + X_{i+1}}{2},\, Y_j,\, N_k - \tfrac{h_v}{2}\right), \\
	\text{Viga Y del piso } k &: \; b_v \times (Y_{j+1} - Y_j - b_c) \times h_v, \quad
	\left(X_i,\, \tfrac{Y_j + Y_{j+1}}{2},\, N_k - \tfrac{h_v}{2}\right).
\end{align}

Una viga X solo se genera si están seleccionadas sus dos intersecciones y el tramo $(i,j) \rightarrow (i{+}1,j)$; igual
para Y. Los nombres son \ruta{Z-}, \ruta{C-piso-}, \ruta{VX-piso-} y \ruta{VY-piso-} seguidos de la marca. La cota
superior de las vigas coincide con el nivel y sus caras extremas tocan las caras de las columnas, de modo que el
metrado y el análisis sísmico las reconocen como unidas.

El tope de 2000 elementos se comprueba antes de generar, con la cota superior

\begin{equation}
	n_{\max} = n_p \left( [\text{zapatas}] + (m) \left( [\text{columnas}] + [\text{vigas X}] + [\text{vigas Y}] \right) \right) \le 2000,
\end{equation}

\noindent donde $n_p$ es el número de intersecciones elegidas, $m$ el número de pisos y $[\cdot]$ vale 1 si la opción
está marcada y 0 si no.

\subsubsection{Superposiciones y vista previa}

\ruta{walterQuispesayhuaPrevisualizar} compara cada elemento nuevo con los existentes y con los demás nuevos. Dos
prismas $A$ y $B$ se consideran superpuestos si se solapan estrictamente en los tres ejes:

\begin{equation}
	\lvert x_A - x_B \rvert < \tfrac{l_A + l_B}{2} - \varepsilon, \quad
	\lvert y_A - y_B \rvert < \tfrac{a_A + a_B}{2} - \varepsilon, \quad
	\lvert z_A - z_B \rvert < \tfrac{h_A + h_B}{2} - \varepsilon, \qquad \varepsilon = 10^{-9}\ \text{m}.
	\label{eq:mod-superposicion}
\end{equation}

Las superposiciones \textbf{no impiden} crear los elementos: se listan en la vista previa para que usted decida.

\subsubsection{Flujo del diálogo}

\begin{enumerate}
	\item El diálogo carga la cuadrícula guardada del proyecto o, si no hay, valores por defecto (separaciones X
	      «5; 5», Y «4», niveles «0; 3»; zapata $1{,}5 \times 1{,}5 \times 0{,}5$\,m, columna de 0,4\,m, viga
	      $0{,}3 \times 0{,}5$\,m). Los valores se separan con punto y coma y admiten coma decimal.
	\item «Actualizar cuadrícula» valida la definición, crea una casilla por intersección y por tramo (todas marcadas) y
	      dibuja la planta con los elementos existentes en gris.
	\item «Revisar elementos» exige que la cuadrícula dibujada coincida con la escrita, genera la vista previa (planta
	      con los nuevos en azul, lista con medidas y centros, y superposiciones) y pide confirmación.
	\item «Crear elementos» llama a \ruta{walterQuispesayhuaCrear}, que guarda la cuadrícula
	      (\ruta{CuadriculaProyecto}, con las listas unidas por «; ») y luego cada elemento en la misma transacción.
	\item «Guardar ejes y niveles» guarda solo la cuadrícula. Los ejes guardados sirven luego como referencias al mover
	      elementos (sección~\ref{sec:mod-movimiento}).
\end{enumerate}

% ------------------------------------------------------------------------------
\subsection{Movimiento en 3D y vínculos viga--columna}
\label{sec:mod-movimiento}

El botón «Mover en 3D» abre \ruta{ui/EditorMovimiento}, que usa \ruta{service/MovimientoElementoService}. El servicio
trabaja con instantáneas inmutables (\ruta{Estado}: id, proyecto, tipo, nombre, medidas y \ruta{Posicion}) en lugar de
entidades, para poder comprobar que nada cambió entre la vista previa y el guardado.

\subsubsection{Concurrencia optimista}

Antes de cualquier operación, \ruta{walterQuispesayhuaComprobarActual(Estado)} relee el elemento de la base y exige
que su estado actual sea \textbf{exactamente igual} al de la instantánea; si no, lanza
\ruta{IllegalStateException} («El elemento cambió desde que comenzó la edición»). Lo mismo se aplica a cada viga
arrastrada por el movimiento y, al deshacer, a cada elemento que el movimiento dejó guardado.

\subsubsection{Sugerencias de alineación}

Si el proyecto tiene una cuadrícula guardada, \ruta{walterQuispesayhuaSugerencias} ofrece, como botones opcionales,
alinear el centro con el eje X o Y más cercano y alinear un nivel. Para columnas se alinea la base
($z - h/2$); para los demás tipos, la cara superior ($z + h/2$). Solo se sugiere una referencia $r$ que cumpla

\begin{equation}
	10^{-6} < \lvert r - v \rvert \le 0{,}25\ \text{m},
\end{equation}

\noindent donde $v$ es el valor actual; entre varias, la más cercana.

\subsubsection{Advertencias}

\ruta{walterQuispesayhuaAdvertencias} produce dos clases de aviso, que no bloquean el guardado:

\begin{itemize}
	\item \textbf{Superposición}: la posición propuesta cumple la ecuación~\eqref{eq:mod-superposicion} con otro
	      elemento (aquí $\varepsilon = 10^{-6}$\,m).
	\item \textbf{Separación}: el elemento movido es una columna que en su posición original estaba en contacto con una
	      viga o zapata y en la propuesta ya no. El contacto admite tocarse:
	      $\lvert x_A - x_B \rvert \le \tfrac{l_A + l_B}{2} + \varepsilon$ en los tres ejes.
\end{itemize}

\subsubsection{Vigas vinculadas a una columna}

Al mover una \textbf{columna}, \ruta{walterQuispesayhuaCandidatos} lista las vigas que ya tienen un vínculo guardado
(\ruta{VinculoGeometrico}, pares columna--viga) o que tocan la columna por uno de sus extremos. El extremo conectado
lo calcula \ruta{walterQuispesayhuaExtremoConectado}. Sea $\tau = 10^{-6}$\,m, sea la viga de dirección X si
$l > a$ (si $l = a$ no se considera viga), $c$ y $v$ las coordenadas de columna y viga a lo largo del eje de la viga,
$w_c$ el lado de la columna en ese eje y $L$ la longitud de la viga. Entonces:

\begin{align}
	\text{condiciones previas} &: \; \lvert z_c - z_v \rvert \le \tfrac{h_c + h_v}{2} + \tau
	\quad \text{y} \quad \lvert p_c - p_v \rvert \le \tau , \\
	e &=
	\begin{cases}
		+1 & \text{si } \left\lvert \left(c + \tfrac{w_c}{2}\right) - \left(v - \tfrac{L}{2}\right) \right\rvert \le \tau ,\\[2pt]
		-1 & \text{si } \left\lvert \left(c - \tfrac{w_c}{2}\right) - \left(v + \tfrac{L}{2}\right) \right\rvert \le \tau ,\\[2pt]
		0 & \text{en otro caso,}
	\end{cases}
\end{align}

\noindent donde $p$ es la coordenada perpendicular (Y para una viga en X). Con $e = +1$ la columna está en el inicio de
la viga (el extremo de menor coordenada) y con $e = -1$, en el final. Es decir, el eje de la viga debe pasar por el
centro de la columna y su cara extrema debe coincidir con la cara de la columna.

El usuario marca qué vínculos quiere guardar. Para cada viga marcada, \ruta{walterQuispesayhuaAjustarViga} mantiene fijo
el extremo lejano y lleva el otro a la cara de la columna en su nueva posición $c'$:

\begin{align}
	f &= v + e \, \tfrac{L}{2}, \qquad m = c' + e \, \tfrac{w_c}{2}, \\
	L' &= e \, (f - m), \qquad v' = \tfrac{f + m}{2},
\end{align}

\noindent con $e = \pm 1$ el extremo conectado. Si $L' \le 10^{-6}$\,m, el movimiento se rechaza («dejaría sin
longitud a la viga»). Una viga vinculada solo puede reajustarse si la columna se mueve \textbf{sobre el eje de la
viga}: $z$ no puede cambiar y tampoco la coordenada perpendicular.

\subsubsection{Operación, guardado y deshacer}

\ruta{walterQuispesayhuaPreparar} arma un registro \ruta{Operacion} con la vista previa, las listas de estados
anteriores y propuestos (el elemento y sus vigas ajustadas) y los conjuntos de vínculos antes y después; añade avisos
de superposición de las vigas reajustadas.

\begin{itemize}
	\item \ruta{walterQuispesayhuaGuardar(Operacion)} vuelve a preparar la operación con los mismos datos y exige que
	      el resultado sea igual al recibido; si no, lanza «El modelo cambió. Revise nuevamente la vista previa». Luego
	      aplica cada estado propuesto y reemplaza los vínculos de la columna.
	\item \ruta{walterQuispesayhuaDeshacer(Operacion)} exige que los vínculos actuales sean los que dejó la operación y
	      que cada elemento siga como quedó; entonces restaura los estados anteriores y los vínculos previos. Si algo
	      cambió, lanza «Los vínculos cambiaron; no se puede deshacer este movimiento» o el error de concurrencia.
\end{itemize}

\subsubsection{Interacción en el editor}

El panel muestra campos X, Y, Z (centro, en metros, con coma decimal), un selector de eje y las casillas de vínculo.
Al arrastrar con el botón principal el elemento que se mueve, el visor convierte el desplazamiento horizontal del ratón
$\Delta u$ (píxeles) en metros a lo largo del eje elegido:

\begin{equation}
	\Delta = \Delta u \cdot \frac{D}{1200} \cdot \frac{1}{65},
\end{equation}

\noindent donde $D$ es la distancia actual de la cámara y 65 la escala del visor en píxeles por metro. Cada cambio en
los campos, las casillas o el arrastre recalcula la operación y la muestra en el visor (incluidas las vigas
reajustadas). «Revisar y guardar» muestra el centro anterior y el propuesto, el volumen, los avisos y los cambios de
cada viga, y guarda tras la confirmación. Mientras dura el movimiento, el controlador deshabilita el resto de la
edición.

% ------------------------------------------------------------------------------
\subsection{Edición en la pantalla del modelo}
\label{sec:mod-pantalla}

\ruta{controller/ModeloController} reúne la lógica de interacción de la pantalla «Modelo 3D y metrado».

\subsubsection{Filtro por tipo y selección múltiple}

La tabla usa una \ruta{FilteredList} sobre todos los elementos del proyecto y una \ruta{SortedList} para ordenar por
columnas; la selección es múltiple. El combo de filtro ofrece las opciones de \ruta{YEFRI\_LERMA\_FILTROS}: «Todos los
tipos», «Solo zapatas», «Solo columnas», «Solo vigas», «Solo losas» y «Solo muros». Cambiar el filtro limpia la
selección. «Todos» selecciona todas las filas visibles. La etiqueta de selección muestra «$N$ en la lista» o «$N$
seleccionados».

\subsubsection{Edición de uno o varios elementos}

«Editar» con una fila abre el panel «Editar elemento» (tipo, nombre, medidas y centro); con varias, el panel «Editar
$N$ elementos». En el segundo, cada medida (largo, ancho, altura) se interpreta con
\ruta{yefriLermaLeerAjuste(String)}:

\begin{table}[H]
	\centering
	\caption{Interpretación de cada campo del panel de lote.}
	\label{tab:mod-lote}
	\begin{tabular}{L{3.4cm}L{4.4cm}L{5.8cm}}
		\hline
		\textbf{Texto} & \textbf{Ajuste} & \textbf{Nueva medida} \\ \hline
		vacío & ninguno (\ruta{null}) & se conserva la actual \\
		\ruta{0,50} & fijo & $m' = 0{,}50$ \\
		\ruta{+0,05} & relativo & $m' = m + 0{,}05$ \\
		\ruta{-0,05} & relativo & $m' = m - 0{,}05$ \\
		\hline
	\end{tabular}
\end{table}

\ruta{yefriLermaAplicarAjustes} devuelve copias (no toca los originales), redondea cada medida a $10^{-4}$\,m
($m' \leftarrow \operatorname{round}(10^4 m') / 10^4$) y rechaza el lote entero si alguna medida queda nula o
negativa. Debe llenarse al menos un campo. El guardado usa \ruta{yefriLermaGuardarVarios}, de modo que el lote se
guarda completo o no se guarda.

Desde el análisis sísmico, \ruta{yefriLermaPrecargarCorreccion} abre este mismo panel con las columnas recomendadas:
precarga una medida fija si todas quedan iguales y, si no, el mayor de los incrementos (por ejemplo «+0,05»), en cada
caso para el largo y el ancho. No guarda nada hasta que usted pulse «Aplicar».

\subsubsection{Aviso de vigas que se sueltan}
\label{sec:mod-sueltas}

Antes de guardar una edición (de uno o de varios elementos), \ruta{yefriLermaConfirmarSiSueltaVigas} llama a
\ruta{yefriLermaVigasQueSeSueltan(actuales, cambiados)}. Este método evalúa
\ruta{FernandoHuancaAnalisisSismico.\allowbreak{}fernandoHuancaConexionVigas} sobre el modelo actual y sobre el modelo con los
cambios aplicados, y devuelve las vigas que \textbf{antes} unían al menos dos columnas y \textbf{después} no. Si la lista
no está vacía se pide confirmación («Guardar igual»); el cambio se puede deshacer. Si el proyecto no tiene columnas o
vigas (el análisis lanza \ruta{IllegalArgumentException}), no hay aviso.

El criterio de conexión es el del análisis sísmico, con tolerancia $T = 0{,}05$\,m. Los pisos son las cotas superiores
de las vigas agrupadas (valores a menos de $T$ del primero del grupo se promedian) que quedan por encima de la base de
las columnas. Una viga de dirección X (si $l \ge a$), con centro $(x_v, y_v)$, longitud $L$ y cota superior en el piso
$s$, «toca» la columna $c$ si

\begin{gather}
	\lvert y_c - y_v \rvert \le \tfrac{a_c}{2} + T, \qquad
	x_v - \tfrac{L}{2} - \tfrac{l_c}{2} - T \le x_c \le x_v + \tfrac{L}{2} + \tfrac{l_c}{2} + T, \\
	z^{-}_c + T < z_s \le z^{+}_c + T,
\end{gather}

\noindent y análogamente en Y. Una viga que toca menos de dos columnas se informa como «suelta».

\subsubsection{Historial de deshacer}

El controlador mantiene una pila (\ruta{ArrayDeque}) de pasos de tipo sellado \ruta{YefriLermaPaso}, con un máximo de
\ruta{YEFRI\_LERMA\_MAXIMO\_DESHACER} $= 20$ pasos (al pasar del máximo se descarta el más antiguo):

\begin{itemize}
	\item \ruta{YefriLermaPasoEdicion}: copias de los elementos antes y después de editar uno o varios.
	\item \ruta{YefriLermaPasoMovimiento}: la \ruta{Operacion} de un movimiento en 3D.
\end{itemize}

«Deshacer» y el atajo \textbf{Ctrl+Z} (acelerador de la escena registrado en \ruta{initialize}) sacan el paso más
reciente. El atajo solo actúa si la tabla está en la escena (la pantalla se está viendo), si no hay un movimiento en
curso y si no hay un panel o diálogo abierto (\ruta{Notificador.\allowbreak{}fernandoHuancaHayVentanaAbierta()}).

\begin{itemize}
	\item Para una edición, \ruta{yefriLermaDeshacerEdicion} compara cada elemento guardado con el estado «después» del
	      paso (\ruta{yefriLermaMismoEstado}: tipo y nombre iguales, medidas y coordenadas con diferencia menor que
	      $10^{-9}$). Si alguno ya no existe o cambió por otro camino, \textbf{rechaza} el deshacer completo; si no,
	      vuelve a guardar los estados «antes» con \ruta{yefriLermaGuardarVarios}.
	\item Para un movimiento, delega en \ruta{MovimientoElementoService.\allowbreak{}walterQuispesayhuaDeshacer}.
	\item El paso se retira de la pila antes de intentarlo: si se rechaza, se muestra el motivo y el paso
	      \textbf{se descarta} (la prueba \ruta{noDeshaceSiElElementoCambioDespuesPorOtroCamino} lo verifica).
\end{itemize}

El historial se vacía al cambiar de proyecto, al eliminar el proyecto y al volver a la pantalla si algún paso es de
otro proyecto. Crear elementos (formulario, plantilla, cuadrícula o DXF) y eliminarlos \textbf{no} se registra en el
historial.

\subsubsection{Clics en el visor y en la tabla}

\begin{table}[H]
	\centering
	\caption{Acciones del ratón sobre la selección (\ruta{yefriLermaClicEn3d}, \ruta{yefriLermaEditarDesde3d} y \ruta{yefriLermaDobleClicFila}).}
	\label{tab:mod-clics}
	\begin{tabular}{L{4.2cm}L{9.4cm}}
		\hline
		\textbf{Acción} & \textbf{Efecto} \\ \hline
		Clic en 3D & Selecciona solo ese elemento. Si ya era parte de una selección múltiple, espera 480\,ms
		(450\,ms de doble clic más 30\,ms) antes de hacerlo, por si llega un segundo clic. \\
		Ctrl+clic en 3D & Agrega o quita el elemento de la selección; si el filtro lo oculta, el filtro pasa a «Todos los
		tipos». También cuentan Shift y Cmd. \\
		Doble clic en 3D & Si el elemento es parte de una selección múltiple, abre «Editar $N$ elementos» con todos; si
		no, lo selecciona y abre «Editar elemento», que incluye el botón «Mover en 3D». \\
		Doble clic en una fila & Igual que en 3D. Como el primer clic deja seleccionada solo esa fila, el controlador
		recuerda la selección previa en el evento \ruta{MOUSE\_PRESSED}. \\
		\hline
	\end{tabular}
\end{table}

\subsubsection{Globo del metrado}

\ruta{yefriLermaActualizarMetrado} pone en el botón el texto «Metrado $\cdot$ $V$ m$^3$», con el volumen neto a dos
decimales, y prepara el resumen completo con \ruta{yefriLermaResumenMetrado}. «Metrado» abre o cierra un
\ruta{Popover} de AtlantaFX anclado al botón con ese resumen; el globo también se cierra cada vez que el metrado se
recalcula, para no mostrar cifras viejas.
«Cuantías de acero\dots» abre el diálogo descrito en la sección~\ref{sec:mod-metrado}.

% ------------------------------------------------------------------------------
\subsection{Visor 3D}
\label{sec:mod-visor}

\ruta{ui/ErickDelgadoVisor3D} dibuja cada elemento como un \ruta{Box} de JavaFX dentro de una \ruta{SubScene} con
antialiasing, cámara en perspectiva (planos de recorte 1 y 30\,000), una luz ambiente y tres puntuales. JavaFX tiene el
eje Y hacia abajo, por lo que el visor transforma las coordenadas del modelo así:

\begin{align}
	(X, Y, Z)_{\text{JavaFX}} &= k \, (x,\; -z,\; y), \qquad k = 65\ \text{px/m}, \\
	(\text{ancho}, \text{alto}, \text{profundidad})_{\text{Box}} &= k\,(l,\; h,\; a).
\end{align}

\begin{table}[H]
	\centering
	\caption{Colores del visor.}
	\label{tab:mod-colores}
	\begin{tabular}{L{4cm}L{3.2cm}L{6.4cm}}
		\hline
		\textbf{Elemento} & \textbf{Color} & \textbf{Observación} \\ \hline
		Zapata & \ruta{\#78716c} & \\
		Columna & \ruta{\#3b82f6} & \\
		Viga & \ruta{\#14b8a6} & también cualquier tipo no previsto \\
		Losa & \ruta{\#94a3b8} & translúcida (opacidad 0,35) \\
		Muro & \ruta{\#f59e0b} & \\
		Seleccionado & \ruta{\#f59e0b} & con brillo especular mayor \\
		Ejes X, Y, Z & \ruta{\#ef4444}, \ruta{\#22c55e}, \ruta{\#3b82f6} & \\
		\hline
	\end{tabular}
\end{table}

\subsubsection{Navegación y vistas}

\begin{itemize}
	\item \textbf{Girar}: arrastre con el botón principal; el ángulo vertical cambia $0{,}5^\circ$ por píxel y queda
	      entre $-10^\circ$ y $90^\circ$.
	\item \textbf{Desplazar}: arrastre con el botón secundario o central, o con Shift; el desplazamiento se escala por
	      $D/1200$.
	\item \textbf{Acercar}: la rueda multiplica la distancia $D$ por 0,9 o la divide entre 0,9, dentro de
	      $[150,\ 9000]$.
	\item \textbf{Vistas} (animadas en 350\,ms por el camino angular más corto): Isométrica ($25^\circ$, $-35^\circ$),
	      Planta ($90^\circ$, $0^\circ$), Frontal ($0^\circ$, $0^\circ$) y Lateral ($0^\circ$, $-90^\circ$).
	      «Ajustar» vuelve a la isométrica y reencuadra. «Cuadrícula» y «Ejes» muestran u ocultan el suelo y los ejes.
\end{itemize}

Al cargar datos (no al seleccionar), \ruta{erickDelgadoConstruirReferencias} dibuja un suelo con líneas cada
$p = \max\!\left(1, \lceil \max(\Delta x, \Delta y) / 40 \rceil\right)$ metros, con 2\,m de margen, situado en la cota
más baja del modelo, y centra el modelo. La distancia de encuadre supone un campo de visión vertical de $30^\circ$:

\begin{equation}
	D = \max\!\left(500,\; \frac{r}{\sin(15^\circ)} \cdot 1{,}1\right), \qquad
	r = \tfrac{1}{2} \sqrt{w^2 + d^2 + \left(k \, H\right)^2},
\end{equation}

\noindent donde $w$ y $d$ son el ancho y el fondo del suelo en píxeles y $H$ la altura total del modelo.

\subsubsection{Selección y modo de movimiento}

\ruta{erickDelgadoMostrar(lista, resaltado, alClic, alDobleClic, reencuadrar)} redibuja las cajas, resalta las que
cumplen el predicado (selección múltiple) e instala los manejadores. Un clic se ignora si hay un movimiento en curso,
si el ratón se movió desde que se presionó o si no es el botón principal. El doble clic se detecta \textbf{por id del
elemento y por tiempo} (450\,ms, \ruta{ERICK\_DELGADO\_DOBLE\_CLIC\_MS}), no por nodo, porque las cajas se recrean al
cambiar la selección; un clic con Ctrl nunca inicia un doble clic. El callback \ruta{AlClic} recibe el elemento y si se
pidió agregar.

\ruta{erickDelgadoIniciarMovimiento(id, alArrastrar)} activa el arrastre del elemento indicado,
\ruta{erickDelgadoPrevisualizar} mueve o redimensiona cajas sin tocar la base de datos y
\ruta{erickDelgadoMostrarReferencias} dibuja en naranja los ejes de la cuadrícula en cada nivel (necesita al menos dos
ejes por dirección).

% ------------------------------------------------------------------------------
\subsection{Metrado neto}
\label{sec:mod-metrado}

\ruta{service/MetradoNeto} calcula el concreto \textbf{sin doble conteo}: donde dos prismas se cruzan, el volumen común
se asigna una sola vez al elemento de mayor prioridad.

\subsubsection{Orden de prioridad}

Los elementos se ordenan de forma estable por tipo: zapata, columna, viga, muro, losa y, al final, cualquier tipo no
previsto. A igual tipo gana el que aparece antes en la lista (el orden de la lista no afecta a los tipos distintos).

\subsubsection{Volumen neto de cada elemento}

Sea $B_1, \dots, B_n$ la secuencia de cajas en ese orden. El volumen neto del elemento $k$ es su volumen menos la
parte ya ocupada por los anteriores:

\begin{equation}
	V_{\text{neto}}(k) = \max\!\left(0,\; V(B_k) - V\!\left( \bigcup_{j<k} \left( B_k \cap B_j \right) \right)\right),
	\label{eq:mod-neto}
\end{equation}

\noindent y la intersección de dos cajas es otra caja, $[\max(x^-_A, x^-_B),\, \min(x^+_A, x^+_B)]$ en cada eje (nula
si algún intervalo queda vacío). La suma de los netos es el volumen de la unión de todos los prismas, y el volumen
«descontado» es $V_{\text{bruto}} - V_{\text{neto}}$.

\subsubsection{Unión de cajas por compresión de coordenadas}

El volumen de la unión de la ecuación~\eqref{eq:mod-neto} se calcula de forma exacta en
\ruta{yefriLermaVolumenUnion}. Se ordenan sin repetir todas las coordenadas límite de las cajas
($X_0 < X_1 < \dots$, igual en Y y Z), lo que divide el espacio en celdas. Cada celda se cuenta una vez si su centro
cae estrictamente dentro de alguna caja:

\begin{gather}
	V\!\left(\textstyle\bigcup_m C_m\right) = \sum_{i}\sum_{j}\sum_{k} (X_{i+1}-X_i)(Y_{j+1}-Y_j)(Z_{k+1}-Z_k)\;
	\mathbf{1}\!\left[ \exists\, m : (\bar{x}_i, \bar{y}_j, \bar{z}_k) \in \operatorname{int} C_m \right], \\
	\bar{x}_i = \tfrac{X_i + X_{i+1}}{2}, \qquad \bar{y}_j = \tfrac{Y_j + Y_{j+1}}{2}, \qquad \bar{z}_k = \tfrac{Z_k + Z_{k+1}}{2}.
\end{gather}

Con una sola caja se devuelve su volumen directamente. Con $m$ cajas hay hasta $(2m-1)^3$ celdas y cada una se compara
con hasta $m$ cajas; como solo intervienen las intersecciones de un elemento con los anteriores, $m$ suele ser pequeño.

\subsubsection{Encofrado y acero}

\begin{table}[H]
	\centering
	\caption{Superficie de encofrado por tipo (\ruta{yefriLermaSuperficieEncofrado}).}
	\label{tab:mod-encofrado}
	\begin{tabular}{L{2.6cm}L{4.6cm}L{6.4cm}}
		\hline
		\textbf{Tipo} & \textbf{Fórmula} & \textbf{Caras consideradas} \\ \hline
		Columna, zapata & $2(l + a)\,h$ & perímetro por altura \\
		Viga & $\max(l,a)\,\left(\min(l,a) + 2h\right)$ & fondo y dos costados \\
		Muro & $2\max(l,a)\,h$ & dos caras \\
		Losa & $l\,a$ & fondo \\
		Otro & 0 & \\
		\hline
	\end{tabular}
\end{table}

El encofrado se calcula con las medidas brutas, \textbf{sin descontar} los contactos entre elementos. El acero se
estima por tipo con una cuantía en kg por m$^3$ de concreto aplicada al volumen \textbf{neto}:

\begin{equation}
	A_t = q_t \sum_{k \in t} V_{\text{neto}}(k) .
\end{equation}

\begin{table}[H]
	\centering
	\caption{Cuantías referenciales de \ruta{yefriLermaCuantiasReferenciales} (kg/m$^3$).}
	\label{tab:mod-cuantias}
	\begin{tabular}{L{3cm}c}
		\hline
		\textbf{Tipo} & \textbf{$q_t$} \\ \hline
		Zapata & 60 \\
		Columna & 150 \\
		Viga & 120 \\
		Muro & 70 \\
		Losa & 80 \\
		Otro & 0 \\
		\hline
	\end{tabular}
\end{table}

\ruta{service/CuantiasAcero} permite cambiarlas. Se guardan en \ruta{cuantias-acero.properties} dentro de la carpeta de
datos (\ruta{RespaldoDatos.\allowbreak{}erickDelgadoCarpetaDatos()}). Al guardar, cada valor debe ser finito y estar entre 0 y
\ruta{YEFRI\_LERMA\_MAXIMO} $= 400$\,kg/m$^3$; si no, no se escribe nada. Al leer, cualquier valor ausente, ilegible o
fuera de rango se reemplaza por el referencial de su tipo, y un archivo inexistente equivale a los referenciales. El
resumen indica «cuantías típicas» o «cuantías propias» según coincidan o no con las referenciales.

\subsubsection{Ejemplo verificado}

La prueba \ruta{MetradoNetoTest.\allowbreak{}descuentaLosCrucesUnaSolaVezSegunPrioridad} usa una columna, una viga que nace en ella y
una losa que cubre ambas, pasadas en orden inverso a la prioridad.

\begin{table}[H]
	\centering
	\caption{Metrado neto del ejemplo de prueba (m$^3$).}
	\label{tab:mod-ejemplo-metrado}
	\begin{tabular}{L{2.2cm}L{5.4cm}ccL{2.6cm}}
		\hline
		\textbf{Elemento} & \textbf{Centro; medidas (m)} & \textbf{Bruto} & \textbf{Neto} & \textbf{Descuento} \\ \hline
		Columna & $(0;\,0;\,1{,}5)$; $0{,}4 \times 0{,}4 \times 3{,}0$ & 0,480 & 0,480 & ninguno \\
		Viga & $(1{,}0;\,0;\,2{,}75)$; $2{,}4 \times 0{,}3 \times 0{,}5$ & 0,360 & 0,300 & cruce con la columna: 0,060 \\
		Losa & $(1{,}0;\,1{,}0;\,2{,}9)$; $4{,}0 \times 4{,}0 \times 0{,}2$ & 3,200 & 3,048 &
		$0{,}032 + 0{,}144 - 0{,}024 = 0{,}152$ \\
		\hline
		\textbf{Total} & & 4,040 & 3,828 & 0,212 \\
		\hline
	\end{tabular}
\end{table}

El término $-0{,}024$ es el cruce triple columna--viga--losa, que la unión cuenta una sola vez.

% ------------------------------------------------------------------------------
\subsection{Importación de planos DXF y DWG}
\label{sec:mod-dxf}

\subsubsection{Lector DXF}

\ruta{service/FernandoHuancaPlanoDxf.\allowbreak{}fernandoHuancaLeer(Path)} lee DXF de texto (ASCII) por pares
código--valor:

\begin{itemize}
	\item Rechaza archivos de más de 10\,MB y archivos truncados, sin \ruta{EOF} o sin entidades.
	\item Elige la codificación por \ruta{\$ACADVER}: UTF-8 desde \ruta{AC1021} (AutoCAD 2007) y Windows-1252 en las
	      versiones anteriores, para no corromper nombres de capa con tildes o eñes. Sustituye los saltos de línea
	      \ruta{CR CR LF} que puede emitir LibreDWG.
	\item De la sección \ruta{HEADER} solo toma \ruta{\$INSUNITS} (código 70). De \ruta{ENTITIES} toma cada entidad con su
	      capa (código 8). Para cada código guarda el \textbf{primer} valor, salvo los vértices de
	      \ruta{LWPOLYLINE} (10, 20) y sus abultamientos (42), que se acumulan.
	\item Una entidad del espacio papel (código 67 igual a 1, o código 410 distinto de \ruta{Model}) se marca como
	      \ruta{PRESENTACIÓN} y se ignora. Una entidad con coordenadas ausentes o no finitas se inventaría como
	      «\textit{tipo} inválida».
\end{itemize}

\begin{table}[H]
	\centering
	\caption{Entidades reconocidas por \ruta{fernandoHuancaEntidad}.}
	\label{tab:mod-dxf-entidades}
	\begin{tabular}{L{3cm}L{3cm}L{7.6cm}}
		\hline
		\textbf{DXF} & \textbf{Tipo interno} & \textbf{Interpretación} \\ \hline
		\ruta{LINE} & \ruta{LINE} & Inicio (10, 20, 30) y fin (11, 21, 31). \\
		\ruta{POINT} & \ruta{POINT} & Centro; requiere un lado indicado por el usuario. \\
		\ruta{INSERT} & \ruta{BLOQUE} & Punto de inserción como centro; requiere un lado. El bloque no se expande. \\
		\ruta{CIRCLE} & \ruta{CÍRCULO} & Cuadrado circunscrito de lado $2r$. \\
		\ruta{LWPOLYLINE} de 2 vértices & \ruta{LINE} & Abierta y sin abultamientos; $z$ = elevación (código 38). \\
		\ruta{LWPOLYLINE} de 4 vértices & \ruta{RECTÁNGULO} & Cerrada (bit 1 del código 70, o 5 vértices con el
		último igual al primero), sin abultamientos, normal $(0,0,1)$, vértices distintos y lados paralelos a X e Y. \\
		Otra \ruta{LWPOLYLINE} & \ruta{LWPOLYLINE} & No compatible con capas estructurales. \\
		\ruta{3DFACE} en capa \ruta{CL3D\_\textit{TIPO}\_\textit{n}} & --- & Cara de un prisma exportado por
		ConcretoLab (ver abajo). \\
		\hline
	\end{tabular}
\end{table}

Si el archivo contiene caras \ruta{3DFACE} en capas \ruta{CL3D\_(COLUMNA|VIGA|ZAPATA|LOSA|MURO)\_\textit{n}}, el lector
reconstruye el modelo 3D completo: para cada capa exige 8 vértices distintos y 6 caras que contengan las del prisma
envolvente; si no, rechaza el archivo («Prisma incompleto o deformado»). El elemento recuperado tiene centro y medidas
del prisma envolvente y el tipo del nombre de la capa.

\subsubsection{Unidades}

\begin{table}[H]
	\centering
	\caption{Conversión de \ruta{\$INSUNITS} (\ruta{fernandoHuancaMetrosPorUnidad}).}
	\label{tab:mod-dxf-unidades}
	\begin{tabular}{cL{3.6cm}c}
		\hline
		\textbf{Código} & \textbf{Unidad} & \textbf{m por unidad} \\ \hline
		1 & pulgadas & 0,0254 \\
		2 & pies & 0,3048 \\
		4 & milímetros & 0,001 \\
		5 & centímetros & 0,01 \\
		6 & metros & 1 \\
		14 & decímetros & 0,1 \\
		otro o ausente & sin unidad declarada & --- \\
		\hline
	\end{tabular}
\end{table}

Si no hay unidad declarada, \ruta{DeteccionCapasDxf.\allowbreak{}erickDelgadoUnidad} la estima por la mayor extensión $E$ del
dibujo (puntos, bloques, rectángulos, círculos y líneas): milímetros si $E > 2000$, centímetros si $E > 200$ y metros en
otro caso. El valor sugerido se coloca en el campo «Metros por unidad DXF», que el usuario puede cambiar. Para un modelo
\ruta{CL3D} recuperado se usa la unidad declarada o, si falta, se asumen metros.

\subsubsection{Detección de capas}

\ruta{service/DeteccionCapasDxf} sugiere una capa por rol a partir de su nombre. \ruta{erickDelgadoPalabras} quita las
tildes (normalización NFD), pasa a mayúsculas y divide el nombre en palabras por cualquier carácter no alfanumérico.
\ruta{erickDelgadoPuntaje} asigna a cada capa el mejor puntaje entre sus palabras $p$ y las claves $q$ del rol:

\begin{equation}
	\text{puntaje} =
	\begin{cases}
		3 & p = q,\\
		2 & \lvert q \rvert \ge 4 \text{ y } \left( p \text{ empieza por } q \text{, o } q \text{ empieza por } p \text{ con } \lvert p \rvert \ge 4 \right),\\
		1 & \lvert q \rvert \ge 5 \text{ y } p \text{ contiene } q,\\
		0 & \text{en otro caso.}
	\end{cases}
\end{equation}

\begin{longtable}{L{2.2cm}L{7cm}L{4.4cm}}
	\caption{Roles, palabras clave y entidades admitidas.} \label{tab:mod-dxf-roles} \\
	\hline
	\textbf{Rol} & \textbf{Palabras clave} & \textbf{Entidades de la capa} \\ \hline
	\endfirsthead
	\hline
	\textbf{Rol} & \textbf{Palabras clave} & \textbf{Entidades de la capa} \\ \hline
	\endhead
	\hline
	\endfoot
	Columnas & COLUMNA, COLUMNAS, COL, COLS, COLUMN, COLUMNS, PILAR, PILARES & puntos, bloques, rectángulos,
	círculos \\
	Vigas & VIGA, VIGAS, VG, BEAM, BEAMS, VIG & líneas \\
	Zapatas & ZAPATA, ZAPATAS, ZAP, CIMENTACION, CIMIENTO, CIMIENTOS, FOOTING, FOOTINGS, FOUNDATION & puntos, bloques,
	rectángulos, círculos \\
	Losas & LOSA, LOSAS, ALIGERADO, ALIGERADOS, SLAB, SLABS, TECHO & rectángulos \\
	Muros & MURO, MUROS, PLACA, PLACAS, WALL, WALLS, MUR, ALBANILERIA & líneas \\
	Ejes & EJE, EJES, GRID, GRIDS, AXIS, AXES, TRAZO & líneas (solo diagnóstico) \\
\end{longtable}

Solo compiten las capas \textbf{compatibles}: con 1 a 200 entidades, todas de los tipos admitidos. La capa de zapatas
no puede ser la misma que la de otro rol ya elegido. \ruta{erickDelgadoDiagnosticar} arma el inventario que ve el
usuario: cada capa con el conteo de entidades por tipo y el rol detectado, más avisos (entidades ignoradas, capas que
parecen de losas o muros pero no son importables, capas con más de 200 entidades, entidades inválidas, círculos
convertidos en secciones cuadradas, bloques usados como punto, espacio papel y unidad estimada).

\subsubsection{Conversión de la planta en edificio}

\ruta{fernandoHuancaConvertir(Plano, Datos)} repite una planta en altura. Todas las coordenadas se multiplican por la
escala $s$ (m por unidad) y deben estar en $z = 0$.

\begin{enumerate}
	\item \textbf{Datos}: de 1 a 12 pisos, altura de piso $h$ mayor que el peralte de viga $h_v$, y escala, altura,
	      ancho y peralte de viga y espesor de zapata positivos.
	\item \textbf{Huellas de columna} (\ruta{fernandoHuancaHuellas}): un rectángulo o círculo da su centro y medidas; un
	      punto o bloque da un cuadrado del lado indicado. Se exigen al menos 4 columnas, sin centros duplicados y a más
	      de 0,05\,m entre sí. Las huellas de zapata (capa opcional) se obtienen igual, con su propio lado.
	\item \textbf{Anclaje de vigas} (\ruta{fernandoHuancaColumna}): cada extremo $p$ de una línea de la capa de vigas
	      se asigna a la única columna cuya huella, ampliada 0,05\,m, lo contiene:
	      \begin{equation}
	      	\lvert p_x - c_x \rvert \le \tfrac{l_c}{2} + 0{,}05, \qquad \lvert p_y - c_y \rvert \le \tfrac{a_c}{2} + 0{,}05 .
	      \end{equation}
	      Si no hay una sola, se toma la de centro más cercano $d_1$ siempre que $d_1 \le 0{,}30$\,m y que la segunda
	      esté al menos a $1{,}5\,d_1$. Si no, se rechaza el plano.
	\item \textbf{Validación de vigas}: extremos en columnas distintas, viga paralela a X o a Y (tolerancia
	      $10^{-6}$), luz libre mayor que el ancho de viga y sin vigas duplicadas. Toda columna debe tener al menos una
	      viga y el conjunto debe ser \textbf{conexo} (se recorre el grafo columna--viga desde la primera columna).
	\item \textbf{Generación} por piso $k = 1, \dots, n$, con $w_a$ y $w_b$ los lados de las columnas extremas a lo
	      largo del eje de la viga:
	      \begin{align}
	      	L &= \lVert b - a \rVert - \tfrac{w_a + w_b}{2}, \\
	      	x_v &= \tfrac{a_x + b_x}{2} + \operatorname{sgn}(b_x - a_x)\, \tfrac{w_a - w_b}{4}
	      	\quad \text{(igual en } y\text{)},
	      \end{align}
	      de modo que la viga va de cara a cara de columna. Las cotas son: zapatas con centro en $-h_z/2$ (cara superior
	      en 0), columnas con centro en $k h - h/2$ y altura $h$, vigas con centro en $k h - h_v/2$.
\end{enumerate}

La variante \ruta{fernandoHuancaConvertir(Plano, Datos, Extra)} agrega en cada piso las \textbf{losas} (rectángulos
de su capa, espesor $e \le h$, centro en $k h - e/2$) y los \textbf{muros} (líneas paralelas a X o Y más largas que su
espesor, de piso a piso, centro en $k h - h/2$). Los muros conservan la longitud de la línea dibujada, sin recortarse
en las columnas; el metrado neto resuelve el cruce. Cada capa estructural admite como máximo 200 entidades y solo los
tipos previstos. Los elementos se nombran \ruta{DXF-\textit{TIPO}-\textit{n}}.

\subsubsection{DWG y RVT}

Para un DWG, \ruta{fernandoHuancaLocalizarDwgread} busca \ruta{dwgread.exe} (LibreDWG) en este orden: la propiedad
\ruta{concretolab.tools}, la variable \ruta{CONCRETOLAB\_TOOLS}, las carpetas \ruta{tools} junto al código de la
aplicación, \ruta{tools} en la carpeta de trabajo, \ruta{\textasciitilde/.concretolab/tools}, \ruta{target} (solo
desarrollo) y el \ruta{PATH}; también revisa subcarpetas cuyo nombre empiece por \ruta{libredwg}.
\ruta{fernandoHuancaLeerDwg} ejecuta \ruta{dwgread -O DXF -o \textit{temporal} \textit{archivo}} con un límite de 45\,s,
lee el DXF temporal y lo borra. Si no encuentra el conversor, la interfaz permite elegirlo. Un archivo RVT no se lee:
se muestran instrucciones para exportarlo a DXF desde Revit.

\subsubsection{Flujo de importación}

\ruta{ui/ErickDelgadoImportacionDxf.\allowbreak{}erickDelgadoAbrir} (menú «+ Proyecto $\rightarrow$ Importar proyecto desde
CAD\dots»):

\begin{enumerate}
	\item Selecciona un archivo \ruta{.dxf}, \ruta{.dxf.txt}, \ruta{.dwg} o \ruta{.rvt} y lo lee (el DWG se convierte en
	      un hilo aparte con un diálogo de espera cancelable).
	\item Si no hay un modelo \ruta{CL3D} y faltan capas compatibles de columnas o de vigas, o si el plano parece vial
	      (capas \ruta{C-ROAD-CORR} y \ruta{C-ROAD-SECT-GRID}) sin capas COLUMNAS y VIGAS, solo muestra el inventario de
	      capas.
	\item Si no, abre «Importar planta estructural» con el inventario, los combos de capa precargados con la
	      sugerencia, la escala sugerida, los siete datos de la planta (pisos, altura de piso, lado de columna, ancho y
	      peralte de viga, lado y espesor de zapata) y las capas opcionales de losas y muros con sus espesores. Los
	      campos que no aplican quedan deshabilitados.
	\item «Revisar modelo» genera los candidatos sin guardar y muestra el conteo por tipo y el concreto neto y bruto
	      (\ruta{MetradoNeto.\allowbreak{}yefriLermaCalcular}).
	\item «Guardar en proyecto» llama a \ruta{yefriLermaGuardarPlano} en un proyecto nuevo o en uno existente vacío.
\end{enumerate}

Con el archivo de ejemplo \ruta{planta-estructural-ejemplo.dxf} aparece el botón «Completar valores del ejemplo», que
carga escala 0,001 y los datos 3; 3; 0,40; 0,30; 0,50; 1,50; 0,50.

% ------------------------------------------------------------------------------
\subsection{Exportación}
\label{sec:mod-exportacion}

\subsubsection{DXF}

\ruta{FernandoHuancaPlanoDxf.\allowbreak{}fernandoHuancaExportar} escribe un DXF ASCII (cabecera \ruta{AC1009}, sin
\ruta{\$INSUNITS}) con cada elemento como seis caras \ruta{3DFACE} en la capa \ruta{CL3D\_\textit{TIPO}\_\textit{n}}.
El color (código 62) es 5 para columnas, 4 para vigas, 30 para muros y 8 para el resto. Las coordenadas se escriben en
metros. El archivo no puede superar 10\,MB, que es el límite del importador, y al reimportarlo se recupera el modelo
exacto (prueba \ruta{PlanoDxfRobustezTest.\allowbreak{}exportarEImportarConservaElModeloExacto}).

\subsubsection{IFC}

\ruta{service/ExportadorIfc.\allowbreak{}erickDelgadoGenerar(nombre, elementos)} genera un archivo IFC2X3 en texto STEP
(ISO 10303-21, \ruta{FILE\_DESCRIPTION} con \ruta{ViewDefinition [CoordinationView]}). La estructura espacial es
proyecto $\rightarrow$ sitio $\rightarrow$ edificio $\rightarrow$ un único nivel «Nivel 0» en la cota 0, unidos con
\ruta{IFCRELAGGREGATES}. Las unidades son SI (metro, m$^2$, m$^3$, radián) y hay un material «Concreto» asociado a
todos los elementos con \ruta{IFCRELASSOCIATESMATERIAL}.

\begin{table}[H]
	\centering
	\caption{Correspondencia de tipos en IFC.}
	\label{tab:mod-ifc}
	\begin{tabular}{L{2.6cm}L{6cm}}
		\hline
		\textbf{Tipo} & \textbf{Entidad IFC2X3} \\ \hline
		Columna & \ruta{IFCCOLUMN} \\
		Viga & \ruta{IFCBEAM} \\
		Muro & \ruta{IFCWALLSTANDARDCASE} \\
		Losa & \ruta{IFCSLAB} (\ruta{.FLOOR.}) \\
		Zapata & \ruta{IFCFOOTING} (\ruta{.PAD\_FOOTING.}) \\
		Otro & \ruta{IFCBUILDINGELEMENTPROXY} \\
		\hline
	\end{tabular}
\end{table}

La geometría de cada elemento es un \ruta{IFCEXTRUDEDAREASOLID}: un perfil rectangular de $l \times a$ centrado en el
origen local, extruido en $+Z$ una longitud $h$, colocado en $(x, y, z - h/2)$ respecto del nivel. Las vigas también
se extruyen en vertical: la forma es exacta, pero el eje de extrusión no sigue la dirección de la viga. Cada elemento
lleva su nombre en \ruta{Name} y su tipo en \ruta{Description}. Los números se escriben con seis decimales y punto.

Los identificadores globales son \textbf{reproducibles}: \ruta{erickDelgadoGuid(semilla)} toma el UUID basado en
nombre de la semilla (\ruta{UUID.\allowbreak{}nameUUIDFromBytes}, MD5) y lo codifica en los 22 caracteres de IFC con el alfabeto
\ruta{0-9A-Za-z\_\$}. Con $n$ el entero de 128 bits y $A$ el alfabeto:

\begin{equation}
	c_0 = A\!\left[\left\lfloor n / 2^{126} \right\rfloor \bmod 4\right], \qquad
	c_{21-i} = A\!\left[\left\lfloor n / 2^{6i} \right\rfloor \bmod 64\right] \quad (i = 0, \dots, 20).
\end{equation}

Las semillas combinan el nombre del proyecto con un prefijo (\ruta{proyecto|}, \ruta{sitio|}, \ruta{edificio|},
\ruta{nivel|}, \ruta{ra1|} a \ruta{ra3|}, \ruta{rc|}, \ruta{rm|}) y, para los elementos, el índice y el nombre
(\ruta{el|\textit{proyecto}|\textit{i}|\textit{elemento}}). Los textos STEP duplican los apóstrofos y las barras
invertidas y codifican cada carácter fuera de ASCII imprimible como
\ruta{\textbackslash X2\textbackslash\textit{hhhh}\textbackslash X0\textbackslash}.

\subsubsection{ETABS (experimental)}

\ruta{service/ExportadorEtabs.\allowbreak{}erickDelgadoGenerar(nombre, fcMpa, elementos)} escribe un archivo de texto \ruta{.e2k}
(cabecera de versión 9.7.4, unidades kN y m) y devuelve también una lista de avisos que la interfaz muestra al
terminar. La resistencia $f'_c$ se pide al usuario (21\,MPa por defecto).

\begin{enumerate}
	\item \textbf{Pisos}: cotas superiores de vigas y losas redondeadas al milímetro (si no hay, las de columnas). La
	      base es la cota inferior más baja de columnas y muros (si no hay, el primer piso menos 3\,m). Se descartan las
	      cotas que no superan la base en más de 5\,mm. Los pisos se llaman \ruta{PISO1}, \ruta{PISO2}\dots; el más alto
	      es \ruta{MASTERSTORY}.
	\item \textbf{Puntos}: únicos en planta a 1\,mm.
	\item \textbf{Columnas}: una línea por columna en su centro, asignada al piso donde termina, con sección
	      \ruta{COL\textit{l}X\textit{a}} (cm). Las columnas que arrancan en la base reciben empotramiento
	      (\ruta{UX UY UZ RX RY RZ}).
	\item \textbf{Vigas}: cada extremo se ancla al centro de la columna cuya huella, ampliada 5\,cm, lo contiene (la más
	      cercana); si no, se usa el extremo de la viga y se avisa. Sección \ruta{VIG\textit{b}X\textit{h}} (cm).
	\item \textbf{Losas}: paneles de 4 puntos con propiedad \ruta{LOSA\textit{e}} (membrana).
	      \textbf{Muros}: paneles con propiedad \ruta{MURO\textit{e}} (\ruta{ShellThin}) definidos con puntos repetidos,
	      lo que se avisa.
	\item \textbf{Zapatas}: no se exportan; se avisa de cuántas se omitieron.
\end{enumerate}

Una columna, viga, losa o muro cuya cota superior no coincide (a 5\,mm) con un piso reconocido se omite y se avisa.

El material \ruta{CONCRETO} se define con peso unitario 23,5\,kN/m$^3$, coeficiente de Poisson 0,2, coeficiente
térmico $9{,}9 \times 10^{-6}$ y

\begin{equation}
	E_c = 4700 \sqrt{f'_c} \ \text{MPa} \;\rightarrow\; 1000 \cdot 4700 \sqrt{f'_c}\ \text{kN/m}^2, \qquad
	f'_c\ (\text{kN/m}^2) = 1000\, f'_c\ (\text{MPa}).
\end{equation}

\begin{advertencia}[Formato sin verificar]
	El propio código marca este exportador como experimental: no se ha comprobado que ETABS importe el archivo. Revise
	pisos, secciones y apoyos en ETABS antes de analizar.
\end{advertencia}

% ------------------------------------------------------------------------------
\subsection{Pruebas automatizadas}
\label{sec:mod-pruebas}

\begin{longtable}{L{6.6cm}L{7.2cm}}
	\caption{Pruebas que cubren este capítulo (\ruta{src/test/java/com/concretolab}).} \label{tab:mod-pruebas} \\
	\hline
	\textbf{Clase} & \textbf{Qué verifica} \\ \hline
	\endfirsthead
	\hline
	\textbf{Clase} & \textbf{Qué verifica} \\ \hline
	\endhead
	\hline
	\endfoot
	\ruta{MetradoNetoTest} & Descuento de cruces una sola vez según prioridad (tabla~\ref{tab:mod-ejemplo-metrado});
	elementos idénticos contados una vez; encofrado y acero de elementos separados. \\
	\ruta{DxfDeteccionTest} & Bloques, círculos y polilíneas con vigas ancladas; rechazo de una viga que no cae en una
	columna; sugerencia de capas y unidad estimada; nombres peruanos con tildes y abreviaturas; losas y muros por piso
	sin duplicar cruces en el metrado. \\
	\ruta{ImportacionDxfTest} & Secciones distintas leídas de los contornos; ida y vuelta de un modelo completo; pórtico
	de dos vanos importado y analizado; planta en milímetros que el motor sísmico puede analizar. \\
	\ruta{PlanoDxfRobustezTest} & Archivos dañados con mensaje claro; límite de 10\,MB; tildes y eñes según la versión;
	conversión de unidades; escalado de un modelo en milímetros; exportar e importar sin pérdida; ubicación del
	conversor DWG. \\
	\ruta{ExportacionIfcEtabsTest} & Estructura IFC, referencias resueltas, GUID de 22 caracteres, codificación de la
	ñ y GUID reproducibles; en ETABS, pisos, secciones, apoyos, avisos y validación de entradas. \\
	\ruta{CuadriculaIntegracionTest} & Generación por ejes en un proyecto existente con detección de superposiciones;
	rechazo de niveles y vanos inválidos sin guardar nada. \\
	\ruta{MovimientoElementoIntegracionTest} & Vista previa con sugerencias y avisos sin guardar; guardar y deshacer
	solo si el elemento no cambió; vínculo y reajuste de vigas (extremo inicial y final) y deshacer de la operación
	completa. \\
	\ruta{ModeloSeleccionMultipleTest} & Edición de varios elementos desde la tabla y desde el visor (clic, Ctrl+clic,
	doble clic); deshacer de varios pasos en orden inverso; rechazo cuando el elemento cambió por otro camino; aviso de
	vigas sueltas con posibilidad de deshacer. \\
	\ruta{controller/EdicionLoteTest} & Lectura de campos vacíos, fijos e incrementos; aplicación solo de los campos
	llenos; precarga de la recomendación sísmica; descripción de la selección. \\
	\ruta{ui/VisorSeleccionTest} & Distinción entre clic, Ctrl+clic y doble clic en el visor y resaltado de la
	selección múltiple. \\
\end{longtable}

% ------------------------------------------------------------------------------
\subsection{Supuestos y límites}
\label{sec:mod-limites}

\begin{itemize}
	\item \textbf{Geometría}: todos los elementos son prismas alineados con los ejes globales. No hay elementos inclinados,
	      circulares ni de sección variable; un círculo del DXF se convierte en una sección cuadrada de lado igual al
	      diámetro.
	\item \textbf{Dirección de las vigas}: no se guarda; se deduce de las medidas, con criterios que no coinciden del
	      todo. El análisis sísmico y el exportador ETABS consideran viga en X si $l \ge a$; el movimiento con vínculos
	      exige $l > a$ (o $a > l$) y no trata como viga una de planta cuadrada; el encofrado usa $\max(l, a)$.
	\item \textbf{Tolerancias}: la generación por cuadrícula usa $10^{-9}$\,m; el movimiento, $10^{-6}$\,m (y 0,25\,m
	      para sugerir alineaciones); el criterio de vigas conectadas del análisis sísmico, 0,05\,m; la importación DXF,
	      0,05\,m de holgura y 0,30\,m de radio de búsqueda; el exportador ETABS, 5\,mm en cotas y 5\,cm en huellas.
	\item \textbf{Metrado}: es una estimación de anteproyecto. El encofrado no descuenta contactos ni distingue caras
	      que no requieren molde; el acero es volumen neto por cuantía fija por tipo. La prioridad de asignación de los
	      cruces es fija (zapata, columna, viga, muro, losa). La pantalla «Mezclas y ensayos» suma volúmenes
	      \textbf{brutos} al elegir elementos del modelo; su propia interfaz advierte que no descuenta intersecciones.
	\item \textbf{Deshacer}: solo cubre ediciones y movimientos (20 pasos); no cubre crear ni eliminar elementos,
	      plantillas, generación por ejes ni importaciones, y un paso rechazado se pierde.
	\item \textbf{Importación DXF}: solo DXF de texto, plantas en $z = 0$ que se repiten idénticas en todos los pisos (1
	      a 12), un máximo de 200 entidades por capa estructural y 10\,MB por archivo. Los bloques no se expanden y no
	      se leen ejes, textos ni cotas para generar elementos. El DWG depende de LibreDWG externo.
	\item \textbf{IFC}: todos los elementos cuelgan de un único nivel en la cota 0 con coordenadas absolutas; no hay
	      armadura ni propiedades (\ruta{Pset}). Los GUID de los elementos dependen de su posición en la lista, por lo que
	      reordenar o borrar elementos cambia los GUID siguientes. No se ha comprobado el archivo con un validador IFC;
	      en particular, \ruta{IFCWALLSTANDARDCASE} se exporta con un material simple y sin eje, cuando el esquema IFC2X3
	      prevé para ese tipo un conjunto de capas de material (verifíquese contra la especificación).
	\item \textbf{ETABS}: experimental y sin verificar; las zapatas se reemplazan por empotramientos y los muros usan
	      paneles con puntos repetidos.
	\item \textbf{Visor}: los muros y la selección usan el mismo color (\ruta{\#f59e0b}), por lo que un muro
	      seleccionado solo se distingue por el brillo especular.
	\item \textbf{Escenarios de laboratorio}: el método existe y está probado, pero no se encontró ninguna llamada desde
	      la interfaz.
\end{itemize}
